\paragraph*{04.05.03.01 Lieferobjekte definieren}
\kcilabel{para:04.05.03.01_LieferobjekteDefinieren}{04.05.03.01}
\label{para:04.05.03.01_LieferobjekteDefinieren}

% Task
\par Nach der Entscheidung zur Anbindung der \gls{adn} an \gls{s4} im Rahmen des \gls{f2h}-Projekts war es meine Aufgabe, die Lieferobjekte für diese Anbindung zu definieren. Dazu gehörte die Identifikation notwendiger Komponenten, Schnittstellen und Datenflüsse sowie die Erarbeitung eines Zielbildes für die Gesamtintegration.

% Action
\par Ich führte eine detaillierte Analyse der bestehenden Systeme und Prozesse durch und arbeitete eng mit den technischen Teams zusammen, um die Schnittstellen und Datenflüsse zu identifizieren (vgl. \autoref{fig:ADNSSTFI}). Mit den Stakeholdern erstellte ich ein umfassendes Dokument, das die Lieferobjekte, ihre Funktionen und Beziehungen zueinander beschreibt.

% Result
\par Das Dokument diente als Grundlage für die weitere Planung und Umsetzung der \gls{adn}-Anbindung. Es sicherte ein gemeinsames Verständnis der Anforderungen und Ziele ab und erleichterte die Kommunikation und Zusammenarbeit zwischen den beteiligten Teams während der Implementierungsphase.

% Reflect
\par Die strukturierte Definition der Lieferobjekte sicherte eine effiziente und nachvollziehbare Umsetzung der \gls{adn}-Anbindung. Dies trug wesentlich dazu bei, dass das Projekt im Zeitplan und Budget geblieben ist. Diese Vorgehensweise werde ich als Best Practice in zukünftigen Wellen des \gls{s4}-Programms anwenden.